\documentclass[aspectratio=169]{beamer}
\usetheme{Madrid}
\usecolortheme{default}
\usepackage[utf8]{inputenc}
\usepackage[T5]{fontenc}
\usepackage[vietnamese]{babel}
\usepackage{array}
\usepackage{booktabs}
\usepackage{hyperref}
\usepackage{tikz}
\usetikzlibrary{positioning,arrows.meta}

\title{Báo cáo giữa kì}
\subtitle{Course Guard:Hệ thống Quản lý Khóa học \& Thi Online}
\author{Nhóm 9}
\institute{NT106 - Lập trình Mạng Căn Bản}
\date{\today}

\begin{document}

%------------------------------------------------
\begin{frame}
  \titlepage
\end{frame}

%------------------------------------------------
\begin{frame}{Mục lục}
  \tableofcontents
\end{frame}

%================================================
\section{4 Tính năng chính hiện tại}

%------------------------------------------------
\begin{frame}{1) Đăng nhập (Login)}
\begin{itemize}
  \item Người dùng nhập \texttt{username/password} tại màn hình Login.
  \item Hệ thống kiểm tra thông tin với dữ liệu trong bảng \texttt{USERS}.
  \item Mật khẩu được so khớp theo cơ chế hash (SHA-256).
  \item Sau khi đăng nhập thành công:
  \begin{itemize}
    \item phân quyền theo role (\texttt{ADMIN/TEACHER/STUDENT}),
    \item ghi log hoạt động (\texttt{AUDIT\_LOGS}),
    \item cập nhật hoạt động thiết bị (\texttt{DEVICES}).
  \end{itemize}
\end{itemize}
\end{frame}

%------------------------------------------------
\begin{frame}{Login - Plant Diagram}
\centering
\begin{tikzpicture}[
  node distance=6mm,
  box/.style={draw, rounded corners, align=center, minimum width=0.72\linewidth, minimum height=8mm},
  line/.style={-{Latex[length=2.2mm]}, thick}
]
  \node[box] (a) {User nhập username/password};
  \node[box, below=of a] (b) {AuthController.Login kiểm tra USERS};
  \node[box, below=of b] (c) {So khớp password hash (SHA-256)};
  \node[box, below=of c] (d) {Đúng: vào Dashboard + ghi AUDIT\_LOGS/DEVICES};
  \node[box, below=of d] (e) {Sai: thông báo đăng nhập thất bại};
  \draw[line] (a) -- (b);
  \draw[line] (b) -- (c);
  \draw[line] (c) -- (d);
  \draw[line] (c.east) -- ++(1.1,0) |- (e.east);
\end{tikzpicture}
\end{frame}

%------------------------------------------------
\begin{frame}{2) Đăng xuất (Logout)}
\begin{itemize}
  \item Người dùng thực hiện Logout từ dashboard.
  \item Hệ thống ghi nhận sự kiện đăng xuất vào \texttt{AUDIT\_LOGS}.
  \item Phiên làm việc kết thúc và điều hướng về màn hình đăng nhập.
  \item Mục tiêu: theo dõi lịch sử truy cập và đảm bảo tính minh bạch vận hành.
\end{itemize}
\end{frame}

%------------------------------------------------
\begin{frame}{Logout - Plant Diagram}
\centering
\begin{tikzpicture}[
  node distance=7mm,
  box/.style={draw, rounded corners, align=center, minimum width=0.72\linewidth, minimum height=8mm},
  line/.style={-{Latex[length=2.2mm]}, thick}
]
  \node[box] (a) {User bấm Logout};
  \node[box, below=of a] (b) {AuthController.Logout()};
  \node[box, below=of b] (c) {Ghi sự kiện LOGOUT vào AUDIT\_LOGS};
  \node[box, below=of c] (d) {Kết thúc phiên và quay về màn hình Login};
  \draw[line] (a) -- (b);
  \draw[line] (b) -- (c);
  \draw[line] (c) -- (d);
\end{tikzpicture}
\end{frame}

%------------------------------------------------
\begin{frame}{3) Quên mật khẩu (Forgot Password)}
\begin{itemize}
  \item User nhập \texttt{username + email} để gửi yêu cầu.
  \item Nếu hợp lệ, trạng thái tài khoản chuyển sang \texttt{RESET\_REQUEST}.
  \item Admin duyệt yêu cầu tại màn hình User Management.
  \item Hệ thống tạo mật khẩu tạm, gửi qua SMTP email cho user.
  \item Sau khi gửi thành công:
  \begin{itemize}
    \item cập nhật \texttt{password\_hash} mới,
    \item chuyển trạng thái user về \texttt{ACTIVE}.
  \end{itemize}
\end{itemize}
\end{frame}

%------------------------------------------------
\begin{frame}{Forgot Password - Plant Diagram}
\centering
\begin{tikzpicture}[
  node distance=5.5mm,
  box/.style={draw, rounded corners, align=center, minimum width=0.78\linewidth, minimum height=8mm},
  line/.style={-{Latex[length=2.2mm]}, thick}
]
  \node[box] (a) {User gửi quên mật khẩu (username + email)};
  \node[box, below=of a] (b) {AuthController.ForgotPasswordRequest()};
  \node[box, below=of b] (c) {Status user $\rightarrow$ RESET\_REQUEST};
  \node[box, below=of c] (d) {Admin duyệt RESET\_REQUEST};
  \node[box, below=of d] (e) {Tạo mật khẩu tạm + gửi SMTP email};
  \node[box, below=of e] (f) {Update password\_hash + status ACTIVE};
  \draw[line] (a) -- (b);
  \draw[line] (b) -- (c);
  \draw[line] (c) -- (d);
  \draw[line] (d) -- (e);
  \draw[line] (e) -- (f);
\end{tikzpicture}
\end{frame}

%------------------------------------------------
\begin{frame}{4) Đăng ký (Register)}
\begin{itemize}
  \item User tạo tài khoản với thông tin: username, full name, email, password.
  \item Tài khoản mới được tạo với:
  \begin{itemize}
    \item role mặc định: \texttt{STUDENT},
    \item status mặc định: \texttt{PENDING}.
  \end{itemize}
  \item Admin duyệt để kích hoạt tài khoản (\texttt{ACTIVE}).
  \item Quy trình này giúp kiểm soát tài khoản ảo và tăng bảo mật đầu vào.
\end{itemize}
\end{frame}

%------------------------------------------------
\begin{frame}{Register - Plant Diagram}
\centering
\begin{tikzpicture}[
  node distance=6mm,
  box/.style={draw, rounded corners, align=center, minimum width=0.78\linewidth, minimum height=8mm},
  line/.style={-{Latex[length=2.2mm]}, thick}
]
  \node[box] (a) {User nhập thông tin đăng ký};
  \node[box, below=of a] (b) {AuthController.RegisterRequest()};
  \node[box, below=of b] (c) {Tạo user mới: role STUDENT, status PENDING};
  \node[box, below=of c] (d) {Admin duyệt tài khoản};
  \node[box, below=of d] (e) {Status user $\rightarrow$ ACTIVE};
  \draw[line] (a) -- (b);
  \draw[line] (b) -- (c);
  \draw[line] (c) -- (d);
  \draw[line] (d) -- (e);
\end{tikzpicture}
\end{frame}

%================================================
\section{Database hiện tại}

%------------------------------------------------
\begin{frame}{Database hiện tại của CourseGuard}
\textbf{Nền tảng:} PostgreSQL trên Supabase (cloud)

\vspace{0.5em}
\textbf{Các nhóm bảng chính:}
\begin{itemize}
  \item \textbf{Nhóm user/auth:} \texttt{USERS, ROLES, PERMISSIONS, ROLE\_PERMISSIONS, DEVICES, AUDIT\_LOGS, NOTIFICATIONS}
  \item \textbf{Nhóm khóa học:} \texttt{COURSES, ENROLLMENTS, ONLINE\_SESSIONS, MATERIALS, MESSAGES}
  \item \textbf{Nhóm thi cử:} \texttt{EXAMS, QUESTIONS, EXAM\_QUESTIONS, EXAM\_ATTEMPTS, ANSWERS, VIOLATIONS}
\end{itemize}

\vspace{0.5em}
\textbf{Chỉ mục đáng chú ý:}
\begin{itemize}
  \item \texttt{IDX\_USERS\_STATUS}
  \item \texttt{IDX\_COURSES\_STATUS}
  \item \texttt{IDX\_AUDIT\_LOGS\_ACTION\_CREATED\_AT}
  \item \texttt{IDX\_AUDIT\_LOGS\_USER\_ID}
\end{itemize}
\end{frame}

%================================================
\section{API sử dụng trong project}

%------------------------------------------------
\begin{frame}{Số lượng API sử dụng trong project}
\textbf{Kiến trúc hiện tại:} ứng dụng WinForms, \alert{không expose REST API công khai}.

\vspace{0.5em}
\textbf{API nội bộ (business/data methods) chính:}
\begin{itemize}
  \item AuthController: \textbf{10} hàm
  \item UserController: \textbf{11} hàm
  \item CourseGuardDbContext: \textbf{20+} hàm truy vấn/cập nhật dữ liệu
  \item SmtpEmailService: \textbf{1} hàm gửi mail chính (\texttt{SendEmail})
\end{itemize}

\vspace{0.5em}
\textbf{Dịch vụ bên ngoài đang tích hợp:}
\begin{itemize}
  \item PostgreSQL endpoint trên Supabase (qua Npgsql)
  \item SMTP server để gửi email thật
\end{itemize}
\end{frame}

%------------------------------------------------
\begin{frame}{Các API quan trọng}
\textbf{Authentication \& User flow}
\begin{itemize}
  \item \texttt{Login()}, \texttt{LoginAsync()}
  \item \texttt{RegisterRequest()}
  \item \texttt{ForgotPasswordRequest()}
  \item \texttt{ApproveUserRequest()}
\end{itemize}

\textbf{Database access}
\begin{itemize}
  \item \texttt{GetUserByUsername()}, \texttt{GetUserByUsernameAndEmail()}, \texttt{GetUserById()}
  \item \texttt{InsertUser()}, \texttt{UpdateUserStatus()}, \texttt{UpdateUserPassword()}
  \item \texttt{SearchUsers()}, \texttt{GetUsersByStatus()}
\end{itemize}

\textbf{Mail service}
\begin{itemize}
  \item \texttt{SmtpEmailService.SendEmail()}
\end{itemize}
\end{frame}

%================================================
\section{Database Supabase \& ưu điểm}

%------------------------------------------------
\begin{frame}{Database đang sử dụng: Supabase PostgreSQL}
\textbf{Lý do lựa chọn:}
\begin{itemize}
  \item Cloud database hoạt động \textbf{24/24}.
  \item \textbf{Không phụ thuộc server localhost} của từng máy cá nhân.
  \item Truy cập ổn định từ nhiều client qua Internet.
  \item Quản trị tập trung, backup và mở rộng dễ dàng hơn mô hình local DB.
\end{itemize}

\vspace{0.5em}
\textbf{Giá trị thực tế với CourseGuard:}
\begin{itemize}
  \item Dễ demo và triển khai cho cả nhóm.
  \item Giảm rủi ro “chạy được trên máy em, máy khác không chạy”.
  \item Hỗ trợ tốt cho các tính năng nhiều người dùng (admin/teacher/student).
\end{itemize}
\end{frame}

%================================================
\section{Kết luận}

%------------------------------------------------
\begin{frame}{Kết luận giữa kỳ}
\begin{itemize}
  \item Đã hoàn thiện 4 tính năng cốt lõi: Login, Logout, Forgot Password, Register.
  \item Hệ thống có luồng quản trị rõ ràng (PENDING, RESET\_REQUEST, ACTIVE).
  \item CSDL được tổ chức đầy đủ cho user, khóa học, thi cử và audit.
  \item Nền tảng Supabase PostgreSQL giúp hệ thống chạy ổn định 24/24, không lệ thuộc localhost.
\end{itemize}

\vspace{0.8em}
\centering
\Large Cảm ơn thầy/cô và các bạn đã lắng nghe!
\end{frame}

\end{document}